\corrigesChapitre{4}{Solutions, ions et pH}
\par\noindent\begin{minipage}{\linewidth}
\begin{corrige}{1}{Classer}
\begin{questions}
\item pH 2 : acide ; pH 11 : basique.
\item Neutre dans les conditions annoncées.
\end{questions}
\end{corrige}
\end{minipage}
\par\noindent\begin{minipage}{\linewidth}
\begin{corrige}{2}{L’unité}
\begin{questions}
\item Le pH est sans unité.
\item Un volume, dans l’unité de sa graduation.
\end{questions}
\end{corrige}
\end{minipage}
\par\noindent\begin{minipage}{\linewidth}
\begin{corrige}{3}{Ions et acidité}
\begin{questions}
\item $\ce{H+}$.
\item $\ce{OH-}$.
\end{questions}
\end{corrige}
\end{minipage}
\par\noindent\begin{minipage}{\linewidth}
\begin{corrige}{4}{Du sel dissous}
\begin{questions}
\item $\ce{Na+}$ et $\ce{Cl-}$.
\item Ils sont dispersés à l’échelle microscopique et traversent le filtre ordinaire avec la solution.
\end{questions}
\pcFigureEx{figures/rev-cor-ions-filtre.pdf}{Un filtre ordinaire retient les grains, pas les ions dissous de cette solution.}
\end{corrige}
\end{minipage}
\par\noindent\begin{minipage}{\linewidth}
\begin{corrige}{5}{Un test bleu}
\begin{questions}
\item Présence d’ions $\ce{Cu^{2+}}$ dans les conditions du test.
\item Non : employer un prélèvement séparé.
\end{questions}
\pcFigureEx{figures/rev-cor-tests-ions.pdf}{Prélèvements séparés : chaque réactif sert à une recherche précise ; ne pas cumuler les réactifs dans le même tube.}
\end{corrige}
\end{minipage}
\par\noindent\begin{minipage}{\linewidth}
\begin{corrige}{6}{Une dilution}
\begin{questions}
\item Oui, car 4 reste inférieur à 7.
\item Non, elle a diminué.
\end{questions}
\pcFigureEx{figures/rev-cor-ph-acide.pdf}{Le pH augmente de 2 à 4, mais reste inférieur à 7 ; la solution est encore acide.}
\end{corrige}
\end{minipage}
\par\noindent\begin{minipage}{\linewidth}
\begin{corrige}{7}{Une base diluée}
\begin{questions}
\item Il s’atténue.
\item Non dans ce modèle usuel : le pH tend vers 7 sans rendre la solution acide.
\end{questions}
\pcFigureEx{figures/rev-cor-ph-base.pdf}{Le pH baisse de 12 à 10, mais reste supérieur à 7 ; la solution est encore basique.}
\end{corrige}
\end{minipage}
\par\noindent\begin{minipage}{\linewidth}
\begin{corrige}{8}{Trois inconnues}
\begin{questions}
\item A acide, B neutre, C basique vers 25 °C.
\item Non : la seule présence d’ions ne donne pas le caractère acide ou basique.
\end{questions}
\pcFigureEx{figures/rev-cor-ph-classement.pdf}{Classer par comparaison à 7, à température usuelle.}
\end{corrige}
\end{minipage}
\par\noindent\begin{minipage}{\linewidth}
\begin{corrige}{9}{Un résultat vert}
\begin{questions}
\item Ion fer(II).
\item Deux électrons.
\end{questions}
\pcFigureEx{figures/rev-cor-tests-ions.pdf}{Prélèvements séparés : chaque réactif sert à une recherche précise ; ne pas cumuler les réactifs dans le même tube.}
\end{corrige}
\end{minipage}
\par\noindent\begin{minipage}{\linewidth}
\begin{corrige}{10}{Neutraliser}
\begin{questions}
\item $\ce{H+}$ et $\ce{OH-}$.
\item L’un peut être introduit en excès ; les quantités et les espèces présentes déterminent le résultat.
\end{questions}
\pcFigureEx{figures/rev-cor-acide-base.pdf}{Une unité H\textsuperscript{+} réagit avec une unité HO\textsuperscript{-} dans ce modèle ; un éventuel excès reste dans la solution.}
\end{corrige}
\end{minipage}
\par\noindent\begin{minipage}{\linewidth}
\begin{corrige}{11}{Échantillons comparables}
\begin{questions}
\item Non, la précision de la comparaison colorée ne permet pas cette écriture.
\item Employer un pH-mètre réglé, avec rinçage et protocole adaptés, puis répéter.
\end{questions}
\end{corrige}
\end{minipage}
\par\noindent\begin{minipage}{\linewidth}
\begin{corrige}{12}{Conduction}
\begin{questions}
\item Électrons dans le métal ; ions positifs et négatifs dans la solution.
\item Oui, neutralité globale et mobilité des charges sont compatibles.
\end{questions}
\pcFigureEx{figures/rev-cor-conduction.pdf}{Dans la solution, des ions des deux signes se déplacent ; ce ne sont pas des électrons libres qui traversent la solution comme dans le métal.}
\end{corrige}
\end{minipage}
\par\noindent\begin{minipage}{\linewidth}
\begin{corrige}{13}{Sans moyenne}
\begin{questions}
\item Le pH ne se moyenne pas et ne décrit pas à lui seul les quantités réactives disponibles.
\item Il faut connaître les solutions et les quantités, ou mesurer le pH final dans un protocole autorisé.
\end{questions}
\end{corrige}
\end{minipage}
\par\noindent\begin{minipage}{\linewidth}
\begin{corrige}{14}{L’eau et le climat}
\begin{questions}
\item Non, elle reste basique vers 25 °C.
\item Oui : acidification désigne ici une baisse de pH, même sans passage sous 7.
\end{questions}
\pcFigureEx{figures/rev-cor-ph-ocean.pdf}{Le pH diminue de 8,2 à 8,0 : acidification, tout en restant dans le domaine basique.}
\end{corrige}
\end{minipage}
\par\noindent\begin{minipage}{\linewidth}
\begin{corrige}{15}{Concevoir une analyse}
\begin{questions}
\item Sur l’un : nitrate d’argent, rechercher un précipité blanc. Sur l’autre : soude, rechercher un précipité bleu. Prévoir les protections et témoins.
\item Chaque test renseigne sur l’ion associé ; l’un n’identifie pas automatiquement l’autre.
\end{questions}
\par\smallskip\noindent\textbf{Pour comprendre.}\enspace Prélever deux portions de solution dans deux tubes propres. Ajouter quelques gouttes de nitrate d’argent au premier pour rechercher les chlorures ; utiliser la soude dans l’autre pour rechercher notamment les ions cuivre(II). Comparer avec un témoin adapté et conclure dans les limites du test.\par
\pcFigureEx{figures/rev-cor-tests-ions.pdf}{Prélèvements séparés : chaque réactif sert à une recherche précise ; ne pas cumuler les réactifs dans le même tube.}
\end{corrige}
\end{minipage}
\par\noindent\begin{minipage}{\linewidth}
\phantomsection\label{cor-defi-4}\pcSolutionDefi{4}{Le laboratoire enquête}\begin{questions}
\item A est acide ; B basique.
\item Le mélange peut réagir, modifier les espèces et créer un risque. Il faut des prélèvements séparés et un protocole validé.
\item A est acide et contient des ions chlorure selon le test. B est basique. Le test cuivre négatif ne permet pas d’affirmer une absence absolue sans connaître la sensibilité du test.
\end{questions}
\pcFigureEx{figures/rev-cor-ph-solutions.pdf}{A, de pH 3, est acide ; B, de pH 10, est basique. Les tests se font sur des prélèvements séparés.}
\end{minipage}
\par\noindent\begin{minipage}{\linewidth}
\phantomsection\label{cor-activite-4}\pcSousTitre{Pistes pour l’investigation}
Noter les valeurs de pH avec une précision compatible avec l’outil. Le test de précipitation se fait sur un autre prélèvement. Pour chaque test, écrire le réactif, l’observation et l’espèce associée dans la banque. Un témoin contrôle une cause indépendante de l’échantillon. Une conclusion négative reste limitée par la sensibilité et les conditions.\par\end{minipage}
